\section{Conclusion}
\label{sec:conclusion}

We have formulated the portability of a comparative scheduler ranking,
$R(W, M, L)$, as a benchmark object distinct from both workload
characterization and single-scheduler evaluation, and designed, froze, and
executed a paired, multi-source evaluation protocol -- a mechanism-diverse
fixed 13-policy scheduler panel, a calibrated six-region operating grid, an
explicit metric contract, and an uncertainty/reversal-detection
methodology -- against $9{,}360$ structurally validated configurations, before
interpreting any comparative outcome. The three primary workload sources
are measurably different and strongly separable in their observed
descriptor distributions (\S\ref{sec:workloads}). Comparative scheduler
rankings are portable in aggregate (Kendall's tau-b never below $0.55$
across 18 source-pair $\times$ region conditions) but source-pair-dependent
rather than uniform: two of the three sources agree almost perfectly with
each other, while the third is measurably and consistently less portable
against either, and this gap holds at every operating region tested
(\S\ref{sec:results}). A small ($3.6\%$) but non-zero share of pairwise
comparisons are practically and statistically supported reversals; all
$36$ involve the same policy pair, with BurstGPT on one side in every
case, concentrated at and above the pre-knee operating region rather than
spread evenly across the load curve (\S\ref{sec:reversals}). Exact-order recovery
of the full-corpus ranking requires most or all of the 40-window budget
for two of three sources, while identifying only the single best policy is
achievable from far fewer windows for those same two sources
(\S\ref{sec:sample-complexity}).

Two post-campaign extensions probe portability further. Across
evaluation metrics, which metric is used changes the comparative ranking
in a non-trivial share of conditions (median Kendall's $\tau_b = 0.419$,
top-1 agreement in only $31.9\%$ of conditions) -- a substantially larger
effect than any source-pair disagreement observed under a single fixed
metric (\S\ref{sec:results-cross-metric}). Across alternative
SLO-synthesis conventions, ranking robustness is comparatively high but
still imperfect, so a headline reversal can persist under one reasonable
SLO definition while not persisting under another
(\S\ref{sec:results-slo}). We report only a pilot-scale, engineering
validation of the synthetic-to-real transfer pipeline for RQ3, not a
transfer conclusion (\S\ref{sec:results-rq3-pilot}). A selected physical
vLLM/GPU validation (RQ6) of this benchmark's largest identified reversal
did not reproduce that reversal, while a paired stable-control comparison
retained the same qualitative ordering on real hardware -- a selected-case
simulator-fidelity limitation, not a general estimate of hardware
reversal prevalence (\S\ref{sec:real-system}).

Ranking portability, in this benchmark, is therefore neither universal nor
uniformly fragile: it is real, but conditional on the workload source,
operating region, evaluation metric, and SLO definition under which a
comparison was established, and a scheduler-selection claim should state
those conditions explicitly rather than being read as general. LSSP is
released as a public, independently reproducible artifact
(\S\ref{sec:artifact}) spanning all six research questions at their
differing levels of completeness -- from the frozen primary analysis
through the pilot-scale RQ3 result and the selected RQ6 physical
validation -- so that each can be independently verified and extended.
